\documentclass{article}
\usepackage[utf8]{inputenc}
\usepackage{amsmath,amssymb}
\usepackage[most]{tcolorbox}
\usepackage{geometry}
\geometry{margin=2.5cm}

\newcommand{\step}[1]{\par\noindent\hangindent=3.5em\hangafter=1%
  \makebox[3.5em][l]{\textbf{#1.}}}
\newcommand{\steptag}[1]{\hfill{\small\textsf{[#1]}}}

\newtcolorbox{proofbox}[1]{
  colback=gray!5, colframe=gray!60,
  fonttitle=\bfseries, title={#1},
  breakable, enhanced,
  left=4pt, right=4pt, top=4pt, bottom=4pt,
  before skip=10pt, after skip=10pt
}

\begin{document}

\begin{proofbox}{If B=⊕\_r B\_r (order-unit norm = max\_r) and each factor ...}
\step{1} If B=⊕\_r B\_r (order-unit norm = max\_r) and each factor has an exact-adjoint coboundary splitting with constant K\_r, then B has one with constant max\_r K\_r + 1, independent of the number of summands (off-block components recovered by P\_r f(e\_r,·), with no sum over r); valid for adjoint/block-respecting modules. \steptag{validated} \\[6pt]
\step{1.1} SETUP / FRAME. By GT-hyp, B = direct-sum over r=1..m of unital JB-algebras B\_r (def-jb-algebra), with coordinatewise Jordan product and order-unit norm ||x|| = max\_r ||x\_r||. Each B\_r is the ideal summand U\_\{e\_r\}B for a central idempotent e\_r of B: by GT-central-idempotent-summand a central idempotent e in a unital Jordan algebra makes U\_e B a direct summand and an ideal, so the e\_r are the central units of the summands; e\_r acts as the identity on its own block B\_r (e\_r circ y = y for y in B\_r) and the coordinatewise product makes B\_r an ideal annihilating the other blocks (y in B\_s, s != r => e\_r circ y = 0). Let P\_r: B -> B\_r be the coordinate projection. The cochain module is the ADJOINT module (def-jordan-coboundary): the action is left Jordan multiplication a.m := a circ m, and the Jordan coboundary is (d\textasciicircum{}1 h)(a,b) = a circ h(b) + h(a) circ b - h(a circ b). By GT-hyp (module restriction) all modules are adjoint/block-respecting, so no Peirce-1/2 mixing occurs; arbitrary modules are excluded. By GT-jb-is-order-unit-space each unital JB-algebra (B\_r and B) is a complete order unit space with order unit the identity and order norm equal to the given norm, so by GT-orderunit-norm-formula the norm is ||a|| = inf\{lambda>0: -lambda*1 <= a <= lambda*1\}; by GT-jb-unit-norm-one the identity has norm 1, hence ||e\_r|| = 1. By GT-hyp (E6) each factor carries an exact-adjoint coboundary splitting S\_r: im(d\textasciicircum{}1\_\{B\_r\}) -> C\textasciicircum{}1(B\_r,B\_r) with d\textasciicircum{}1\_\{B\_r\} S\_r g = g and ||S\_r g|| <= K\_r ||g|| in the injective order-unit cochain norm (def-injective-cochain-norm); set K := max\_r K\_r. This node fixes notation and the standing hypotheses used by all later nodes. \steptag{validated} \\[4pt]
\step{1.2} BLOCK RESTRICTION IS A COBOUNDARY. Fix an exact-adjoint coboundary f = d\textasciicircum{}1 h on B (so f in im(d\textasciicircum{}1), with primitive 1-cochain h: B -> B; this f is the input we must split). For each r define the same-factor restriction f\textasciicircum{}r: B\_r x B\_r -> B\_r by f\textasciicircum{}r(a,b) := P\_r f(a,b) for a,b in B\_r, and define the restricted 1-cochain h\_r: B\_r -> B\_r by h\_r(a) := P\_r h(a). CLAIM: f\textasciicircum{}r = d\textasciicircum{}1\_\{B\_r\} h\_r, i.e. f\textasciicircum{}r is an exact-adjoint coboundary on B\_r. PROOF: for a,b in B\_r, unfold the B-coboundary (def-jordan-coboundary, node 1.1) and apply P\_r: P\_r f(a,b) = P\_r[ a circ h(b) + h(a) circ b - h(a circ b) ]. Since B\_r is an ideal of B and a,b in B\_r (node 1.1), a circ b in B\_r, and for any z in B the product a circ (P\_r z) and the projection commute on the B\_r block: P\_r(a circ h(b)) = a circ P\_r h(b) = a circ h\_r(b), because a in B\_r and the coordinatewise product has no cross-block contribution to the B\_r component (node 1.1: B\_r is a block-ideal, e\_r acts as identity on B\_r and annihilates other blocks, so P\_r(a circ w) = a circ P\_r w for a in B\_r). Likewise P\_r(h(a) circ b) = (P\_r h(a)) circ b = h\_r(a) circ b, and P\_r h(a circ b) = h\_r(a circ b). Hence f\textasciicircum{}r(a,b) = a circ h\_r(b) + h\_r(a) circ b - h\_r(a circ b) = (d\textasciicircum{}1\_\{B\_r\} h\_r)(a,b). Therefore f\textasciicircum{}r in im(d\textasciicircum{}1\_\{B\_r\}), so S\_r f\textasciicircum{}r is defined (node 1.1, E6) and d\textasciicircum{}1\_\{B\_r\}(S\_r f\textasciicircum{}r) = f\textasciicircum{}r. \steptag{validated} \\[4pt]
\step{1.3} CONSTRUCTION OF THE SPLITTING Sf. Given the exact-adjoint coboundary f = d\textasciicircum{}1 h on B (node 1.2), define a 1-cochain Sf: B -> B coordinatewise by its blocks: for x in B and each r, (Sf)\_r(x) := P\_r (Sf)(x) is given by (Sf)\_r(x) = (S\_r f\textasciicircum{}r)(x\_r) + P\_r f(e\_r, x\_\{!=r\}), where x\_r := P\_r x is the B\_r-component, x\_\{!=r\} := x - x\_r is the off-block part, S\_r is the per-factor splitting and f\textasciicircum{}r = d\textasciicircum{}1\_\{B\_r\} h\_r the same-factor restriction (node 1.2). Then Sf(x) := sum\_r (Sf)\_r(x) assembling the blocks. The FIRST summand (S\_r f\textasciicircum{}r)(x\_r) is the diagonal term (well-defined since f\textasciicircum{}r in im(d\textasciicircum{}1\_\{B\_r\}) by node 1.2); the SECOND summand P\_r f(e\_r, x\_\{!=r\}) is the off-block recovery term, evaluated at the central unit e\_r (node 1.1) on the FIRST slot and the off-block part on the second, with NO sum over r (each output block B\_r uses only its own e\_r). Sf is linear in f (both summands are linear in f: S\_r is linear and f $|-\rangle$ f\textasciicircum{}r, f $|-\rangle$ f(e\_r,.) are linear) and linear in x. This defines the candidate map S: im(d\textasciicircum{}1) -> C\textasciicircum{}1(B,B), f $|-\rangle$ Sf. \steptag{validated} \\[4pt]
\step{1.4} COORDINATE MAPS ARE NORM-NONINCREASING (max-norm). Working in the order-unit norm of B and the injective order-unit cochain norm (def-injective-cochain-norm, node 1.1). From ||x|| = max\_r ||x\_r|| (GT-hyp E5, node 1.1): (i) for each r, ||x\_r|| <= max\_s ||x\_s|| = ||x||, so the coordinate projection P\_r is norm-nonincreasing, ||P\_r y|| <= ||y|| for all y in B. (ii) ||x\_\{!=r\}|| = ||x - x\_r|| = ||sum\_\{s != r\} x\_s|| = max\_\{s != r\} ||x\_s|| <= max\_s ||x\_s|| = ||x|| (the max-norm of the off-block part is the max over s != r, dominated by the full max). (iii) ||e\_r|| = 1 (node 1.1, via GT-jb-unit-norm-one). CONSEQUENCE for cochain restrictions: for the same-factor restriction f\textasciicircum{}r(a,b) = P\_r f(a,b) with a,b in B\_r (node 1.2), ||f\textasciicircum{}r||\_inj = sup\_\{||a||,||b|| <= 1, a,b in B\_r\} ||P\_r f(a,b)|| <= sup\_\{||a||,||b|| <= 1 in B\} ||f(a,b)|| = ||f||\_inj, using (i) (P\_r norm-1) and that the sup over the B\_r-unit-ball is over a subset of the B-unit-ball (||a|| <= 1, a in B\_r implies ||a|| <= 1 in B since the norm is inherited, node 1.1). Hence ||f\textasciicircum{}r|| <= ||f|| for every r. \steptag{validated} \\[4pt]
\step{1.5} NORM BOUND ||Sf|| <= (max\_r K\_r + 1)||f||. Take x in B with ||x|| <= 1; bound each output block (Sf)\_r(x) = (S\_r f\textasciicircum{}r)(x\_r) + P\_r f(e\_r, x\_\{!=r\}) (node 1.3). DIAGONAL term: ||(S\_r f\textasciicircum{}r)(x\_r)|| <= ||S\_r f\textasciicircum{}r||\_inj * ||x\_r|| <= K\_r ||f\textasciicircum{}r|| * ||x\_r|| <= K\_r ||f|| * 1 <= K ||f||, using ||S\_r g|| <= K\_r ||g|| (node 1.1, E6), ||f\textasciicircum{}r|| <= ||f|| (node 1.4), ||x\_r|| <= ||x|| <= 1 (node 1.4(i)), and K = max\_r K\_r (node 1.1) so K\_r <= K. OFF-BLOCK term: ||P\_r f(e\_r, x\_\{!=r\})|| <= ||P\_r|| * ||f(e\_r, x\_\{!=r\})|| <= ||f||\_inj * ||e\_r|| * ||x\_\{!=r\}|| <= ||f|| * 1 * 1 = ||f||, using ||P\_r y|| <= ||y|| (node 1.4(i)), the definition of the injective cochain norm ||f(u,v)|| <= ||f||\_inj ||u|| ||v|| (def-injective-cochain-norm, node 1.1), ||e\_r|| = 1 (node 1.4(iii)), and ||x\_\{!=r\}|| <= ||x|| <= 1 (node 1.4(ii)). SUM per block: ||(Sf)\_r(x)|| <= K||f|| + ||f|| = (K+1)||f|| for EVERY r, where the +1 is the single flat cost of the one application of f in the off-block term, INDEPENDENT of the factor constants K\_r. AGGREGATE: ||Sf(x)|| = max\_r ||(Sf)\_r(x)|| <= (K+1)||f|| (the aggregation across blocks is a MAX, by ||y|| = max\_r ||y\_r||, node 1.1/1.4 — NOT a sum). Taking sup over ||x|| <= 1: ||Sf||\_inj <= (K+1)||f|| = (max\_r K\_r + 1)||f||. The bound is INDEPENDENT OF THE NUMBER OF SUMMANDS m: the only cross-summand aggregations (||x|| = max\_r ||x\_r|| and ||Sf|| = max\_r ||(Sf)\_r||) are maxima of per-r quantities each bounded by the same (K+1), so no term scales with m. \steptag{validated} \\[4pt]
\step{1.6} OFF-BLOCK RECOVERY IDENTITY. Write the primitive's output-blocks as h\_\{rs\}: B\_s -> B\_r, h\_\{rs\}(y) := P\_r h(y) for y in B\_s (so the relevant unit is e\_r on the OUTPUT side B\_r). CLAIM: for r != s and y in B\_s, P\_r f(e\_r, y) = h\_\{rs\}(y). PROOF: f = d\textasciicircum{}1 h (node 1.2), so unfold (def-jordan-coboundary, node 1.1) at (e\_r, y): f(e\_r, y) = e\_r circ h(y) + h(e\_r) circ y - h(e\_r circ y). Now e\_r circ y = 0 since y in B\_s, s != r and e\_r annihilates other blocks (node 1.1), so the third term vanishes: h(e\_r circ y) = h(0) = 0. Apply P\_r to the remaining two terms. FIRST term: P\_r(e\_r circ h(y)) = e\_r circ P\_r h(y) = e\_r circ h\_\{rs\}(y) = h\_\{rs\}(y), because e\_r acts as the identity on its own block B\_r (node 1.1: e\_r circ z = z for z in B\_r) and P\_r h(y) = h\_\{rs\}(y) lies in B\_r; also e\_r circ (anything) projected to B\_r equals e\_r circ (its B\_r-part) since e\_r is the block-r central unit (node 1.1, GT-central-idempotent-summand). SECOND term: P\_r(h(e\_r) circ y) = 0, because y in B\_s with s != r and the coordinatewise product sends B\_? circ B\_s into B\_s-related blocks with zero B\_r-component when s != r (node 1.1: distinct blocks multiply to 0 in the coordinatewise product, so any product involving the factor y in B\_s has zero component in B\_r for r != s). Hence P\_r f(e\_r, y) = h\_\{rs\}(y) + 0 = h\_\{rs\}(y). This reads the off-block primitive component h\_\{rs\} (r != s) directly off the coboundary f via one evaluation at the central unit e\_r. \steptag{validated} \\[4pt]
\step{1.7} RIGHT-INVERSE: d\textasciicircum{}1 Sf = f. We show the constructed Sf (node 1.3) satisfies (d\textasciicircum{}1 Sf)(a,b) = f(a,b) for all a,b in B, where d\textasciicircum{}1 is the B-coboundary (def-jordan-coboundary, node 1.1) and f = d\textasciicircum{}1 h is the given coboundary (node 1.2). Both d\textasciicircum{}1 Sf and f are SYMMETRIC bilinear in (a,b), and B = direct-sum\_r B\_r, so by bilinearity it suffices to verify the identity on homogeneous pairs (a,b) with a in B\_r and b in B\_s for all r,s. Two cases exhaust these homogeneous pairs: (CASE A) the same-factor pairs r = s, established in child node 1.7.1; (CASE B) the cross-factor pairs r != s, established in child node 1.7.2. Since every (a,b) in B x B expands by bilinearity into a finite sum of such homogeneous pairs and both sides are bilinear, child nodes 1.7.1 and 1.7.2 together give (d\textasciicircum{}1 Sf)(a,b) = f(a,b) on all of B x B; hence d\textasciicircum{}1 Sf = f, i.e. Sf is an exact right inverse on im(d\textasciicircum{}1). SCOPE OF THIS NODE vs THE TRANSITIVE 1.5 DEPENDENCY. The right-inverse property d\textasciicircum{}1 Sf = f established above uses ONLY this node's declared dependencies -- nodes 1.2 (block restriction is a coboundary), 1.3 (construction of Sf), 1.6 (off-block recovery identity) and the case children 1.7.1, 1.7.2 -- and asserts NO norm constant. The full conclusion the contract headlines -- that the linear map S: f $|-\rangle$ Sf (node 1.3) is an exact-adjoint coboundary splitting on B WITH CONSTANT max\_r K\_r + 1 -- is NOT discharged by this node alone: it is obtained by COMBINING this right-inverse (d\textasciicircum{}1 Sf = f) with the norm bound ||Sf|| <= (max\_r K\_r + 1)||f|| proved in node 1.5. Node 1.5 is therefore a TRANSITIVE DEPENDENCY of that combined, constant-bearing conclusion; we RECORD it here in-statement because af provides no post-hoc dependency-edge command to add 1.5 to this node's edges without re-refining (which would archive the already-validated case children 1.7.1, 1.7.2) -- the same in-prose transitive-dependency recording precedent established at thm-faithful-approx node 1.4. The final assembly -- construction S: f $|-\rangle$ Sf (node 1.3) + norm bound ||Sf|| <= (max\_r K\_r + 1)||f|| (node 1.5) + exact right inverse d\textasciicircum{}1 Sf = f (this node 1.7) => S is an exact-adjoint coboundary splitting on B with constant max\_r K\_r + 1, in the adjoint/block-respecting module -- is discharged at the root node 1. \steptag{validated} \\[6pt]
\step{1.7.1} \textsc{case}: CASE A (same-factor, r = s): a, b in B\_r. Compute (d\textasciicircum{}1 Sf)(a,b) = a circ Sf(b) + Sf(a) circ b - Sf(a circ b) (def-jordan-coboundary, node 1.1) and compare to f(a,b) block by block on the output. INPUTS: a,b in B\_r so a circ b in B\_r (block-ideal, node 1.1); the relevant components of Sf at the inputs a,b,a circ b in B\_r are, for each output block j, (Sf)\_j(a) = (S\_j f\textasciicircum{}j)(a\_j) + P\_j f(e\_j, a\_\{!=j\}) with a\_j = a if j=r else 0 (node 1.3, since a in B\_r). OUTPUT BLOCK j = r: only the B\_r-components of a circ Sf(b), Sf(a) circ b, Sf(a circ b) contribute, and since a,b in B\_r and B\_r is a block-ideal (node 1.1), P\_r(a circ Sf(b)) = a circ (Sf)\_r(b), P\_r(Sf(a) circ b) = (Sf)\_r(a) circ b, P\_r Sf(a circ b) = (Sf)\_r(a circ b). For inputs in B\_r the off-block term of (Sf)\_r vanishes (P\_r f(e\_r, a\_\{!=r\}) with a\_\{!=r\} = 0), so (Sf)\_r restricted to B\_r equals (S\_r f\textasciicircum{}r) restricted to B\_r = S\_r f\textasciicircum{}r composed with identity, i.e. (Sf)\_r|\_\{B\_r\} = S\_r f\textasciicircum{}r (a 1-cochain on B\_r). Hence P\_r (d\textasciicircum{}1 Sf)(a,b) = a circ (S\_r f\textasciicircum{}r)(b) + (S\_r f\textasciicircum{}r)(a) circ b - (S\_r f\textasciicircum{}r)(a circ b) = (d\textasciicircum{}1\_\{B\_r\}(S\_r f\textasciicircum{}r))(a,b) = f\textasciicircum{}r(a,b) = P\_r f(a,b), using d\textasciicircum{}1\_\{B\_r\}(S\_r f\textasciicircum{}r) = f\textasciicircum{}r (node 1.2) and f\textasciicircum{}r = P\_r f|\_\{B\_r x B\_r\} (node 1.2). OUTPUT BLOCK j != r: a circ (Sf)\_j(b) has zero B\_j-component (a in B\_r, the product a circ z lands in B\_r-related blocks, zero in B\_j for j != r; node 1.1) and likewise (Sf)\_j(a) circ b = 0 in B\_j; so P\_j(a circ Sf(b)) = 0 = P\_j(Sf(a) circ b). Thus P\_j (d\textasciicircum{}1 Sf)(a,b) = - P\_j Sf(a circ b) = - (Sf)\_j(a circ b). Now a circ b in B\_r and j != r, so (Sf)\_j(a circ b) = (S\_j f\textasciicircum{}j)((a circ b)\_j) + P\_j f(e\_j, (a circ b)\_\{!=j\}); here (a circ b)\_j = 0 (since a circ b in B\_r, j != r), and (a circ b)\_\{!=j\} = a circ b, giving (Sf)\_j(a circ b) = P\_j f(e\_j, a circ b). By the off-block recovery identity (node 1.6) with output-block j, source-block r, and a circ b in B\_r (r != j): P\_j f(e\_j, a circ b) = h\_\{jr\}(a circ b) = P\_j h(a circ b). On the other side, P\_j f(a,b) = P\_j(d\textasciicircum{}1 h)(a,b) = P\_j[a circ h(b) + h(a) circ b - h(a circ b)] = -P\_j h(a circ b) (the first two terms vanish in B\_j: a,b in B\_r, j != r, so a circ h(b) and h(a) circ b have zero B\_j-component by the block-product rule, node 1.1). Hence P\_j (d\textasciicircum{}1 Sf)(a,b) = -(Sf)\_j(a circ b) = -P\_j h(a circ b) = P\_j f(a,b). Both output blocks agree, so (d\textasciicircum{}1 Sf)(a,b) = f(a,b) for a,b in B\_r. \steptag{validated} \\[4pt]
\step{1.7.2} \textsc{case}: CASE B (cross-factor, r != s): a in B\_r, b in B\_s. Since distinct blocks multiply to zero in the coordinatewise product (node 1.1), a circ b = 0, so the third coboundary term vanishes: (d\textasciicircum{}1 Sf)(a,b) = a circ Sf(b) + Sf(a) circ b - Sf(a circ b) = a circ Sf(b) + Sf(a) circ b (def-jordan-coboundary, node 1.1; Sf(0) = 0 by linearity, node 1.3). Compare to f(a,b) = (d\textasciicircum{}1 h)(a,b) = a circ h(b) + h(a) circ b - h(a circ b) = a circ h(b) + h(a) circ b (node 1.2; same vanishing). Check each output block j. OUTPUT BLOCK j = r: P\_r(a circ Sf(b)) = a circ (Sf)\_r(b) (a in B\_r, block-ideal, node 1.1), and P\_r(Sf(a) circ b) = 0 because b in B\_s, s != r, so any product with the factor b in B\_s has zero B\_r-component (node 1.1). Now b in B\_s and (Sf)\_r(b) = (S\_r f\textasciicircum{}r)(b\_r) + P\_r f(e\_r, b\_\{!=r\}); since b in B\_s with s != r, b\_r = 0 (so the diagonal term is S\_r f\textasciicircum{}r(0) = 0) and b\_\{!=r\} = b, giving (Sf)\_r(b) = P\_r f(e\_r, b). By the off-block recovery identity (node 1.6) with output-block r, source-block s, b in B\_s (r != s): P\_r f(e\_r, b) = h\_\{rs\}(b) = P\_r h(b). Hence (Sf)\_r(b) = P\_r h(b), so a circ (Sf)\_r(b) = a circ P\_r h(b) = a circ h\_\{rs\}(b), giving P\_r (d\textasciicircum{}1 Sf)(a,b) = a circ h\_\{rs\}(b). On the other side P\_r f(a,b) = P\_r[a circ h(b) + h(a) circ b] = a circ P\_r h(b) + 0 = a circ h\_\{rs\}(b) (the term h(a) circ b has zero B\_r-component since b in B\_s, s != r, node 1.1; and P\_r(a circ h(b)) = a circ P\_r h(b) since a in B\_r, node 1.1). So P\_r (d\textasciicircum{}1 Sf)(a,b) = a circ h\_\{rs\}(b) = P\_r f(a,b). OUTPUT BLOCK j = s: symmetric (swap roles of a,b and r,s using symmetry of the Jordan product and of the cochains): P\_s(Sf(a) circ b) = b circ (Sf)\_s(a) and (Sf)\_s(a) = P\_s f(e\_s, a) = h\_\{sr\}(a) = P\_s h(a) (node 1.6, output-block s, source-block r, a in B\_r, s != r); P\_s(a circ Sf(b)) = 0 (a in B\_r, r != s, node 1.1). Hence P\_s (d\textasciicircum{}1 Sf)(a,b) = (Sf)\_s(a) circ b = h\_\{sr\}(a) circ b = P\_s h(a) circ b = P\_s f(a,b) (on the other side P\_s[a circ h(b) + h(a) circ b] = 0 + P\_s h(a) circ b = h\_\{sr\}(a) circ b). OUTPUT BLOCK j with j != r and j != s: P\_j(a circ Sf(b)) = 0 (the factor a in B\_r forces zero B\_j-component for j != r, node 1.1) and P\_j(Sf(a) circ b) = 0 (factor b in B\_s forces zero B\_j-component for j != s, node 1.1), so P\_j (d\textasciicircum{}1 Sf)(a,b) = 0; and P\_j f(a,b) = P\_j[a circ h(b) + h(a) circ b] = 0 likewise (node 1.1). All output blocks agree, so (d\textasciicircum{}1 Sf)(a,b) = f(a,b) for a in B\_r, b in B\_s, r != s. \steptag{validated} \\[4pt]
\end{proofbox}

\end{document}
